Randomized controlled trials identify causal effects under transparent design assumptions, but their sample sizes can make average and conditional treatment-effect estimation imprecise. Observational studies are often larger and can yield accurate outcome predictions, but those predictions need not be causal. We develop a model-agnostic prediction-powered framework that keeps the randomized trial as the causal anchor while using observational outcome models and covariates through two tunable channels. The coefficient \(\lambda\) changes the outcome regression inside an RCT-valid pseudo-outcome, and \(\omega\) introduces a prediction-powered control variate centered across the two samples.

For the average treatment effect, we derive exact conditional unbiasedness, an exact two-sample variance formula, the joint population variance minimizer, its population-orthogonal simplification, and an honest plug-in algorithm. For conditional effects, we formulate a prediction-powered squared-loss learner from an RCT-valid AIPW pseudo-outcome, distinguish target preservation from empirical loss-variance tuning, and give implementable linear-sieve and frozen-neural-basis versions. We outline an R-loss analogue only as a prospective extension requiring a separate derivation. Exact density-ratio transport covers unequal covariate marginals, while loss-targeted balancing gives an approximate alternative with explicit drift. Existing simulations support the ATE construction in several small-trial regimes. CATE experiments and real-outcome analyses remain planned, so the present empirical evidence does not establish a universal finite-sample improvement guarantee.
